% 原书物理页=印刷页63--64；4.2全节，止于独立几何引理证毕.
\section{有限维单模的分类}\label{sec:ssl-4-2}
因为所有有限维不可约表示都是最高权表示，所以它们均具有 $V(\lambda)$ 的形式，其中 $\lambda\in\mathfrak h^*$ 唯一. 因此剩下要确定何时 $V(\lambda)$ 有限维. 我们已经见过，一个必要条件是 $\lambda\in P$. 但这并非充分条件：
\begin{theorem}\label{thm:ssl-4-2-1}
设 $\lambda\in\mathfrak h^*$，$V(\lambda)$ 是最高权为 $\lambda$ 的唯一单 $\mathfrak g$-模. 则 $V(\lambda)$ 有限维，当且仅当对一切 $\alpha\in\Delta$，都有 $\lambda(H_\alpha)\in\mathbb N$.
\end{theorem}
\begin{proof}
先作如下观察. 设 $v$ 是 $V(\lambda)$ 的一个本原向量，$\alpha\in\Delta$. 则 $v$ 生成一个最高权为 $\lambda(H_\alpha)$ 的 $\mathfrak{sl}_\alpha$-模 $V(\lambda,\alpha)$.
\begin{enumerate}[label=\roman*)]
\item 若 $V(\lambda)$ 有限维，则 $V(\lambda,\alpha)$ 也有限维，而我们知道有限维 $\mathfrak{sl}_2$-模的最高权是自然数，即 $\lambda(H_\alpha)\in\mathbb N$.
\item 反过来，假设对一切 $\alpha\in\Delta$ 都有 $\lambda(H_\alpha)\in\mathbb N$. 我们的目标是证明 $P(V(\lambda))$ 在 $W$ 的作用下稳定，再应用下面的几何引理，它告诉我们这时 $P(V(\lambda))$ 有限；由于各 $V(\lambda)_\mu$ 都有限维，这就足以推出 $V(\lambda)$ 有限维.
\end{enumerate}
因为 $W$ 由单反射生成，只需证明 $P(V)$ 在各 $s_\alpha$（$\alpha\in\Delta$）下稳定. 注意，若 $\mu\in P(V)$，则由我们的假设以及 $P(V)$ 的形式，有 $\mu(H_\alpha)\in\mathbb Z$. 若 $\mu(H_\alpha)=0$，则 $s_\alpha(\mu)=\mu$，无需再说什么. 若 $\mu(H_\alpha)>0$，则
\begin{align*}
F_\alpha^{\mu(H_\alpha)}(V(\lambda)_\mu)\subset V(\lambda)_{s_\alpha(\mu)},
\end{align*}
因此只需证明 $F_\alpha^{\mu(H_\alpha)}(V(\lambda)_\mu)\neq0$. 若 $\mu(H_\alpha)<0$，则
\begin{align*}
E_\alpha^{-\mu(H_\alpha)}(V(\lambda)_\mu)\subset V(\lambda)_{s_\alpha(\mu)},
\end{align*}
因此只需证明 $E_\alpha^{\mu(H_\alpha)}(V(\lambda)_\mu)\neq0$.\footnote{译注：原文此处指数为 $\mu(H_\alpha)$，与前一式中的 $-\mu(H_\alpha)$ 不同；因本段假设 $\mu(H_\alpha)<0$，此处疑漏负号，按原文照录.}
而我们知道，在 $\mathfrak{sl}_\alpha$ 的有限维表示 $W$ 中，对 $m\in\mathbb N$ 有 $F_\alpha^m(W_m)=W_{-m}$ 以及 $E_\alpha^m(W_{-m})=W_m$. 因此，为了推出 $s_\alpha(P(V))=P(V)$，只需说明 $V$ 是有限维 $\mathfrak{sl}_\alpha$-子模之和.
于是考虑 $V$ 的所有有限维 $\mathfrak{sl}_\alpha$-子模之和 $V(\lambda)^\alpha$，先注意它在 $\mathfrak g$ 下稳定. 事实上，若 $V'$ 是一个有限维 $\mathfrak{sl}_\alpha$-子模，则 $\sum_{X\in\mathfrak g}XV'$ 也是一个有限维 $\mathfrak{sl}_\alpha$-子模（注意，求和可以只遍历 $\mathfrak g$ 的一组基）. 我们想证明 $V(\lambda)^\alpha=V(\lambda)$. 由 $V(\lambda)$ 的不可约性，现在只需证明 $V(\lambda)^\alpha\neq0$，也就是给出一个有限维 $\mathfrak{sl}_\alpha$-子模.
我们将证明，证明开头定义的 $V(\lambda,\alpha)$ 有限维. 因为 $v$ 对 $\mathfrak b_\alpha=KH_\alpha\oplus\mathfrak g_\alpha$ 是特征向量，所以 $V(\lambda,\alpha)=\operatorname{Vect}_K\{F_\alpha^iv,\ i\in\mathbb N\}$. 回顾公式
\begin{align*}
E_\alpha.F_\alpha^iv=i(\lambda(H_\alpha)-i+1)F_\alpha^{i-1}v,
\end{align*}
它已在研究 $\mathfrak{sl}_2$ 的不可约表示时用过（并且很容易用归纳法证明）. 特别地，它表明 $E_\alpha.F_\alpha^{\lambda(H_\alpha)+1}(v)=0$. 另一方面，对于 $\beta\in\Delta\setminus\{\alpha\}$，我们知道 $E_\beta$ 与 $F_\alpha$ 交换，故
\begin{align*}
E_\beta.F_\alpha^{\lambda(H_\alpha)+1}(v)=F_\alpha^{\lambda(H_\alpha)+1}.E_\beta v=0
\end{align*}
（因为 $v$ 是本原向量）. 因此向量 $F_\alpha^{\lambda(H_\alpha)+1}(v)$ 是本原向量. 而它的权为 $\lambda'=\lambda-(\lambda(H_\alpha)+1)\alpha\neq\lambda$. 由定理~\ref{thm:ssl-4-1-6} 的第 \ref{ssl-4-1-6-ii} 项，$\lambda'$ 不是 $V(\lambda)$ 的最高权，所以 $F_\alpha^{\lambda(H_\alpha)+1}(v)=0$. 于是 $V(\lambda,\alpha)$ 有限维.
\end{proof}
\begin{lemma*}
设 $C:=\sum_{\alpha\in\Delta}\mathbb R_+\alpha\in V$ 是根系 $\Phi$ 的一组基 $\Delta$ 生成的正锥.\footnote{译注：原文此处写作 $C\in V$；锥 $C$ 是 $V$ 的子集，此处照录原记号.}
则对任意 $x\in V$，$\bigcap_{w\in W}w(x-C)$ 是紧的.
\end{lemma*}
\begin{proof}
设 $w_0$ 是 $W$ 中满足 $w_0(\Delta)=-\Delta$ 的唯一元素（它的阶为 $2$）. 只需证明 $(x-C)\cap w_0(x-C)$ 是紧的. 设 $(\omega_\alpha)_\alpha$ 是 $V^*$ 中 $\Delta$ 的对偶基. 于是 $w_0(\omega_\alpha)=-\omega_{-w_0(\alpha)}$. 由定义，
\begin{align*}
x-C=\{v\in V,\ \forall\alpha\in\Delta,\ \langle\omega_\alpha,v\rangle\leq\langle\omega_\alpha,x\rangle\}
\end{align*}
且
\begin{align*}
w_0(x-C)&=\{v\in V,\ \forall\alpha\in\Delta,\ \langle\omega_\alpha,w_0(v)\rangle\leq\langle\omega_\alpha,x\rangle\}\\
&=\{v\in V,\ \forall\alpha\in\Delta,\ \langle-\omega_{-w_0(\alpha)},v\rangle\leq\langle\omega_\alpha,x\rangle\},
\end{align*}
因此\footnote{译注：下一式下标中的 $\omega_0$ 原文如此；上文相应记号为 $w_0$，此处照录.}
\begin{align*}
(x-C)\cap w_0(x-C)
&=\{v\in V,\ \forall\alpha\in\Delta,\\
&\qquad\langle-\omega_{-\omega_0(\alpha)},x\rangle\leq\langle\omega_\alpha,v\rangle\leq\langle\omega_\alpha,x\rangle\},
\end{align*}
这个集合显然是紧的，因为 $(\omega_\alpha)_\alpha$ 是 $V^*$ 的一组基.
\end{proof}
% 原书正文至此结束；以上两个proof的闭合均对应原PDF中的证毕方框.
